\section{Recovery as a Multi-Axis Measurement}
\label{sec:observation}

Perplexity aggregates token-level negative log-likelihood over a corpus. A
knowledge-sensitive task instead asks whether a particular answer can be
selected or generated under a task interface. CPT can therefore improve the
former without changing the latter at the same observed rate. This does not make
perplexity uninformative; it means that distributional fit is not a certificate
for every behavior after a structural intervention.

This distinction motivates three study choices. First, we retain within-arm
checkpoints because a single endpoint cannot distinguish delayed from absent
change over the observed budget. Second, answer-letter and complete-option
scoring are reported separately because they mix competence and readout
differently. Third, we compare structural constructions because the same nominal
depth need not imply the same inherited computation. These are measurement
requirements, not a causal theory of where capability resides.

Layer-wise probes provide supplementary context: on intact OLMo-2-7B, frozen
semantic classifiers become readable early, MMLU answer-letter readout rises
near layers 18--19, and exact next-token agreement saturates at the top. Because
these quantities use different probes and thresholds, we treat their ordering
as descriptive only; full results are in Appendix~\ref{app:mechanism}.
